\chapter{Discussion}
\label{chap:discussion}

This chapter steps back from the numbers and asks what the seven-variant ablation and the Stage~2 results actually support. I revisit each research question in turn, then discuss what the frozen-splat regime is doing to the loss landscape, then the position compression I observed in Stage~2 and what I tried to fix it, and finally the limitations I am aware of and that constrain the conclusions one can draw. I write this chapter in the spirit of an honest project post-mortem rather than a victory lap: the architectural changes I propose are large quantitative improvements, and the Stage~2 generator does work, but the system as a whole is not yet at a quality I would call photorealistic, and I want to be precise about where the residual difficulty lies.

%% ──────────────────────────────────────────────────────────────────────────
\section{Revisiting the Research Questions}
\label{sec:disc_rqs}

\paragraph{RQ1 (Ordering).} The result is positive but with one qualifier. Hilbert serialisation in a fixed canonical frame substantially improves validation reconstruction (V1 $\to$ V2: $99 \to 77$ on the full-scene validation $\ell_2$, Table~\ref{tab:7_summary}). The improvement is entirely a generalisation effect; on the training set, the V1 architecture can also bring the loss down, but it does so by storing per-scene context in the bottleneck. Imposing a canonical order on the target makes the position component of the loss into a learnable function, so the network can use the same parameters across scenes instead of memorising a per-scene mapping. The qualifier is that the ordering only buys what it buys; pairing it with a position scaffold gets a much larger drop (V2 $\to$ V3: $77 \to 45$), so the order is the necessary precondition rather than the full solution.

\paragraph{RQ2 (Architecture and generalisation).} The result is also positive but with a clearer mechanism than I expected. Replacing the flat $777$M-parameter decoder with a $1.6$M-parameter token-local MLP \emph{by itself} does not close the gap (V1 $\to$ V4: stays at $99$). Adding a local encoder \emph{by itself} would not be testable cleanly because the local encoder gives a per-token latent that the flat decoder cannot use coherently. The two changes only work together. When both are on (V5), reconstruction improves to $79$ with a $500\times$ reduction in decoder parameters; when the Hilbert order plus scaffold are added on top (V6), reconstruction drops to about $47$, comparable to V3's $45$ but with two orders of magnitude fewer decoder parameters. The mechanism is that the local encoder and the token-local decoder are two halves of one structural decision: latent token $k$ is a code for spatial region $k$, and the decoder MLP only has to model what is in that region. This compositional structure is what generalises. The qualitative figure (Figure~\ref{fig:recon}) shows that the visual jump happens between V3 and V6 even though the aggregate metric only moves a few points, which is consistent with this reading.

\paragraph{RQ3 (Semantic structure).} The 300-scene experiments (\S\ref{sec:results_300}) show that per-Gaussian InfoNCE based on ScanNet-72 labels can be added at no measurable cost to reconstruction, while producing per-Gaussian features that cluster by category (Figure~\ref{fig:pca}). The layout/geometry split (Strategy~B1's deterministic-layout idea, lifted into the VAE form for the large-scale runs) gives a more stable layout code under partial-scan masking. At scale (V6 vs.\ V3) the disentangled and per-Gaussian-supervised variant matches the un-disentangled variant on reconstruction, and provides the schema (a 16-token $z_s$ plus a 512-token $z_g$) that the hierarchical Stage~2 generator needs. The Stage~2 results confirm that this schema is usable: \texttt{LayoutDiT} reaches a low validation loss on the 16-token layout latent and \texttt{GeometryDiT512} conditioned on it reaches a comparably low loss on the 512-token geometry latent. The qualifier here is that semantic structure was not the limiting factor for reconstruction at any scale studied, so the disentanglement is best read as a Stage-2-enabling design choice rather than a Stage-1 reconstruction win.

\paragraph{Stage 2 generation versus completion.} The qualitative Stage~2 result on V6 (Figure~\ref{fig:stage2_exp6}) is honestly mixed. Unconditional generation reproduces the structural pattern of the Stage~1 latent well; the shape of the generated point cloud closely matches the Stage~1 reconstruction of a held-out scene. Completion under $40\%$ observation recovers the coarse layout but is more diffuse than the reference, with the high-density region spreading along the diagonal of the scene. Two interpretations are consistent with this: completion is a constrained inverse problem on an imperfect latent, and the Stage~1 decoder is particularly sensitive to residual error in the position channel, so any inaccuracy in the completed latent shows up as positional smearing in the rendered output. I view the result as evidence that the Stage~2 architecture and training pipeline are sound (the flow matching converges, generation works), and that the residual limitation lies in the precision of the inferred latent for unobserved regions.

%% ──────────────────────────────────────────────────────────────────────────
\section{What the Frozen-Splat Regime Does to the Loss Landscape}
\label{sec:disc_frozen}

The contrast between the Can3Tok original demonstration (a few hundred scenes, splat-fitting jointly optimised with the encoder) and the SceneSplat-7K setting (thousands of scenes, splats pre-fitted and frozen) deserves its own paragraph because it is, in my reading, the largest residual obstacle to closing the photorealism gap. In Can3Tok's joint regime, the splat-fitting process can subtly co-adapt to whatever per-token tokenisation the encoder ends up using. The decoder slot $i$ has implicit influence on what gets put into the input position $i$. In the pre-fitted regime, this co-adaptation is broken: the splats are fixed by a renderer-quality objective that has no knowledge of the autoencoder, so the network must absorb whatever within-region arrangement the renderer produced.

This explains why the same architecture that worked at $300$ scenes in the joint setting fails to generalise at $6{,}000$ in the pre-fitted setting. The 300-scene exploration in this thesis (also pre-fitted) does not contradict this; at 300 scenes the network is still in the memorisation regime where per-scene context is cheap to store, and the failure mode of the pre-fitted setting is not yet active. The seven-variant ablation is in some sense a study of how to make a slot-aligned loss work when the data-side cooperation that the joint setting silently relied on is removed.

It also explains the otherwise surprising finding that the position scaffold (V3) gives a much larger improvement than ordering alone (V2). With ordering alone, the network still has to predict the absolute coordinate at slot $i$, which is hard for the same reason: within a Hilbert block the absolute position varies room by room. With the scaffold, the network only has to predict a bounded local offset relative to a token-specific anchor that is itself constrained to be the block centroid. The hard work is delegated to the anchor head, which has a structurally easy job, and the offset head solves a simple local-deformation problem. The two-stage decomposition matches what the pre-fitted regime can actually identify from data.

%% ──────────────────────────────────────────────────────────────────────────
\section{Position Compression and What I Tried}
\label{sec:disc_compression}

Across all Stage~2 diagnostics, the positional output of the Stage~1 decoder, when fed Stage~2 latents, lies in a smaller range than the canonical $\pm 10$\,m training frame. V6 generates positions in roughly $\pm 2$\,m; V7 generates them in approximately $\pm 0.7$\,m. Real validation scenes from the same set occupy roughly $\pm 5$ to $\pm 9$\,m in the canonical frame. The compression is consistent across epochs, across the layout, geometry, and completion stages, and across the four scenes saved at each evaluation checkpoint.

I diagnosed the source carefully. The flow-matching cosine similarity between predicted and target velocity sits above $0.96$ for V6 geometry from epoch 50 onwards, which means the DiT is learning the velocity field correctly: the latent it produces matches the latent the encoder produces. The compression therefore is not introduced by Stage~2. Encoding a real validation scene and decoding it directly through the frozen Stage~1 decoder reproduces the same compressed range, which places the issue in the Stage~1 decoder output rather than in any Stage~2 module. The anchor-relative position decoding (Eq.~\eqref{eq:anchor_relative}) bounds the local offset to $\rho \tanh(o_i)$ with $\rho$ a learnable global scale; my best reading is that the network has learned a small $\rho$ during Stage~1 training because the slot-aligned $\ell_2$ loss rewards conservative predictions when the anchor is approximately correct, and the bias persists through the frozen decoder used at Stage~2.

What I tried to fix it during the project:
\begin{itemize}
\item \textbf{Loosening the anchor offset bound.} Re-initialising $\rho$ at a larger value (4.0 instead of 2.0) for a partial re-training run did not change the qualitative compression. The network re-converges to a similar effective scale.
\item \textbf{Diagnostic on the decoded coordinates.} I verified that the Stage~1 decoder receives a latent in the expected range during Stage~2 sampling and that the inverse normalisation in the PLY writer is applied symmetrically to GT and generated PLYs, ruling out a writer-side bug.
\item \textbf{Comparing V6 and V7.} The compression is worse in V7 (no scaffold) than V6 (scaffold), which suggests the scaffold is doing the right thing structurally but the anchor centroids themselves are not at the right absolute scale.
\end{itemize}

What I would do next if the project continued: re-train Stage~1 with an explicit positional-range loss that penalises the output for being narrower than the input, or, more cleanly, parameterise $\rho$ per Hilbert block rather than as a global scalar so that high-extent regions of the scene can use a wider offset range. Both are tractable but were outside the time budget. The honest framing is that Stage~2 demonstrates the latent generative pipeline works in principle, while Stage~1 still has a positional-range bug that holds back the final rendered quality.

%% ──────────────────────────────────────────────────────────────────────────
\section{Limitations}
\label{sec:disc_limitations}

I list the limitations I am aware of, ordered roughly from most consequential.

\paragraph{Stage 1 has not converged to a photorealistic level.} Even the best Stage~1 variant (V6, validation $\ell_2 = 47.5$) produces reconstructions that are recognisable as ``a room with multiple distinct objects'' but are far from a per-scene 3DGS fit in visual quality (Figure~\ref{fig:recon}). The aggregate metric and the qualitative figure both point in the same direction. The thesis demonstrates that the architectural changes I propose move the system from a complete-failure regime (V1) to a regime where the scene structure becomes visible (V6), but the system is not yet at a quality that competes with per-scene 3DGS reconstructions or with the original Can3Tok demonstrations at the 300-scene scale.

\paragraph{Position compression in the decoder output.} As described in \S\ref{sec:disc_compression}, the Stage~1 decoder produces positions in a smaller range than the canonical training frame, and this propagates to all Stage~2 outputs. I have characterised the source but the fix was not in scope for this thesis.

\paragraph{No rendering-quality metrics.} Every number in Chapter~\ref{chap:results} is on the same per-primitive $\ell_2$ residual that is used during training (Stage~1) or the rectified-flow velocity MSE (Stage~2). I do not report PSNR, SSIM, or LPIPS on rendered images of the reconstructed scenes. The per-primitive $\ell_2$ is a useful relative metric for comparing variants but is not directly tied to perceptual quality. Adding rendered-image evaluation is the next concrete step.

\paragraph{Stage 2 completion is partial.} Generation works on the V6 latent (Figure~\ref{fig:stage2_exp6}b), reproducing the structural pattern of the Stage~1 latent. Completion (panel d) is recognisably scene-shaped but more diffuse than the reference. The training loss converges and the flow diagnostics are clean, so I attribute the qualitative gap to the constrained-inversion nature of completion combined with the Stage~1 decoder's sensitivity to small latent errors in the position channel.

\paragraph{Fixed-$N$ assumption.} I work with a fixed $40{,}000$ Gaussians per scene, sampled by top-opacity ranking. Real 3DGS scenes have variable counts (often in the hundreds of thousands), and the top-opacity ranking biases the input distribution towards opaque furniture surfaces. A density-uniform farthest-point sampling would give a more spatially balanced input and is a clear next experiment.

\paragraph{Indoor-only.} SceneSplat-7K is indoor only. The Hilbert canonical frame (radius $10$\,m) is sized for indoor extents. Outdoor scenes with unbounded depth would need either a different normalisation (log-depth contracted) or a tiled-scene approach. I have not tested either.

\paragraph{No comparison to a Can3Tok-style joint training at scale.} A fair test of the frozen-vs-joint hypothesis (\S\ref{sec:disc_frozen}) would be to also train a joint-splat variant at the same scale and compare. Doing this is computationally expensive (splat fitting per scene is not free) and was out of scope for this thesis.

\paragraph{Hyperparameter sweep is shallow.} For each of the seven variants I have a single training run on a single seed. Variance estimates would be informative, but each run takes a couple of days on $2 \times $H100 and the budget did not allow for repeats. The trends between variants are large enough relative to plausible seed variance that I am confident in the qualitative ordering, but precise effect sizes should be treated as point estimates.

\paragraph{ScanNet-72 labels do not transfer cleanly across sub-datasets.} I disable per-Gaussian semantic supervision for the ARKitScenes and ScanNet++ portions of the training corpus because their taxonomies do not match ScanNet-72 exactly. This is a pragmatic choice; a unified label space (or a multi-taxonomy head) would let the per-Gaussian supervision act on every training scene.

\paragraph{Disentanglement is structural, not identifiable.} The layout/geometry split in V6/V7 relies on auxiliary heads (scene-mean colour, scene-label distribution, per-category centroids), cross-reconstruction, and the orthogonality regulariser. These are structural inductive biases. As Locatello et al.\ \citep{locatello2019challenging} point out, unsupervised disentanglement is not identifiable; my split is identifiable only insofar as the auxiliary supervision pins down what $z_s$ is supposed to contain. The supervision could be stronger or weaker without my having a principled way to choose.

None of these limitations invalidates the central finding of the thesis (the joint role of Hilbert ordering, anchor-relative decoding, and the local encoder plus token-local decoder pair in scaling the Stage~1 VAE, and the schema-aware Stage~2 generator that follows from it). They do delimit the scope within which that finding is presently supported.